\chapter{Introducción y Objetivos}

\section{Objetivos}
\begin{itemize}
    \item Diferenciar las potencias activa, reactiva y aparente en sistemas monofásicos de corriente alterna.
    \item Determinar la diferencia de fase entre tensión y corriente utilizando el método de figuras de Lissajous
          en el modo X-Y del osciloscopio.
    \item Diseñar y calcular la compensación del factor de potencia mediante el agregado de capacitores en paralelo.
    \item Medir y comparar la respuesta de multímetros de valor medio calibrados y multímetros de valor eficaz
          verdadero (True RMS) frente a señales no sinusoidales.
    \item Determinar los errores sistemáticos y el grado de incertidumbre asociado a las mediciones realizadas.
\end{itemize}

\section{Introducción Teórica}
En circuitos de corriente alterna de baja frecuencia y frecuencia industrial, las cargas reactivas (inductivas o
capacitivas) introducen un desfasaje entre las señales de tensión y corriente. Esto da origen a tres expresiones
distintas de potencia que caracterizan el comportamiento energético del sistema: la potencia activa ($P$), que representa
la energía útil convertida en trabajo; la potencia reactiva ($Q$), asociada al intercambio de energía para el
sostenimiento de campos eléctricos y magnéticos; y la potencia aparente ($S$), que representa el producto directo del
valor eficaz de tensión por la corriente.

El factor de potencia, definido como $\cos\phi$ (donde $\phi$ es el ángulo de desfasaje), es un parámetro fundamental
para evaluar la eficiencia del transporte y distribución de energía eléctrica. Un factor de potencia bajo incrementa la
corriente necesaria para una misma potencia activa, elevando las pérdidas por efecto Joule en las líneas de transmisión.
Por este motivo, se realiza la compensación del factor de potencia mediante el acoplamiento de bancos de capacitores.

Para caracterizar el desfasaje de forma indirecta, el osciloscopio puede configurarse en modo X-Y para graficar una
figura de Lissajous. La elipse resultante permite deducir el desfase mediante la relación de sus intersecciones con los
ejes y su amplitud máxima.

Por otra parte, la gran mayoría de los multímetros digitales comunes para corriente alterna miden el valor medio
rectificado de la señal y aplican un factor de escala de $1,11$ para indicar el valor eficaz bajo la hipótesis de una
forma de onda puramente sinusoidal. Al medir señales con alto contenido armónico, como ondas cuadradas, triangulares o
con control de fase por Triac (dimmer), se introducen errores significativos debido a que la relación entre el valor
eficaz y el valor medio de módulo (factor de forma) difiere de la del caso sinusoidal. En estas situaciones, es
indispensable utilizar instrumentos con detección True RMS (valor eficaz verdadero).
